\documentclass[letterpaper]{article}
\usepackage[submission]{aaai2027}
\usepackage[hyphens]{url}
\usepackage{graphicx}
\urlstyle{rm}
\def\UrlFont{\rm}
\usepackage{natbib}
\usepackage{caption}
\usepackage{booktabs}
\pdfinfoomitdate=1
\pdfinfo{/TemplateVersion (2027.1)}
\frenchspacing
\setcounter{secnumdepth}{0}
\title{When Stale Context Changes Decisions:\\The Importance of Incidental-Mention Controls}
\author{Anonymous authors}
\affiliations{}
\begin{document}
\maketitle
\begin{abstract}
An old value may still sit in a record long after it stops governing the decision. We edited one word at a time in synthetic routing records, keeping the current state, policy, and correct action fixed. With thinking disabled, all five models shifted their action scores and changed their answers when the word appeared as a superseded value. The same happened when the record simply mentioned the word. In some models the mention produced a larger net increase in wrong answers. Because the word pointing to the wrong action also names the other entity's current state, we cannot tell which cue drove an individual error. In a post hoc check, we still saw action changes on histories where separate questions elicited both the initial and current states correctly. For systems that act on updated records, we propose testing what the system does when those old lines change as well as asking it what the record currently says.
\end{abstract}

\section{Introduction}
Unit A is now \texttt{restricted}, so it goes EAST. An earlier line says it was \texttt{pending}, which would send it WEST under the same policy. Replace \texttt{pending} with \texttt{suspended}, which maps to EAST. The answer should still be EAST. Suppose the model changes its answer. Would it do the same if the earlier line had merely mentioned \texttt{pending}?

That is the comparison we make. We edit one word in four settings. The word can appear as an obsolete routing value, an old archive tag, a mention by the unit, or a mention without a named unit. We score both possible actions and record the model's generated answer. We focus on routing because the policy is an explicit lookup. Models struggled with the separate access task even on current-only records, so we report it in the appendix.

\paragraph{Related work.}
STALE documents a gap between retrieving updated evidence and acting on it \citep{chao2026stale}. Controlled In-Context Memory studies how models return obsolete values after updates \citep{guo2026context}. Contextual entrainment experiments show that irrelevant mentions can raise token probabilities and distract models \citep{niu2025entrainment}. Matched paraphrases can induce as many flips as targeted counterfactual edits, and irrelevant context can induce item-level flips that aggregate accuracy conceals \citep{yang2026compared,zhang2026illusion}. We compare alternate roles of the same word while holding the correct downstream action fixed. CICM reports little cross-variable leakage in its parallel-variable task. Our other-attribute controls concern generated policy actions under different record templates and do not establish leakage in direct state retrieval.

\section{Design and measurements}
\begin{table}[t]
\centering
\begin{tabular}{@{}lp{0.61\columnwidth}@{}}
\toprule
Construction & Earlier line containing the edit \\
\midrule
S: superseded & Initially, A had routing state \textbf{[word]}.\\
O: other attribute & Initially, A had archive tag \textbf{[word]}.\\
M: entity mention & Initially, A mentioned the word \textbf{[word]}.\\
U: unattached & Initially, the unassigned word was \textbf{[word]}.\\
\bottomrule
\end{tabular}
\caption{Shortened routing confirmation prompt, with entity names abbreviated. Swap only \textbf{pending} (WEST) and \textbf{suspended} (EAST) within each row. Both members retain A's final routing state \texttt{restricted} (EAST), B's final state \texttt{pending} (WEST), and A's later archive tag \texttt{restricted}. The query asks for A's action. The correct answer is EAST in all eight prompts.}
\label{tab:example}
\end{table}
\begin{table*}[t]
\centering
\begin{tabular}{@{}llrcccc@{}}
\toprule
Model & Layouts & $n$ & S score (nats) & S net & M net & M $-$ S net \\
\midrule
Qwen3-8B & int., seq. & 256 & $3.248$ ${[2.768, 3.751]}$ & 0.195 & 0.312 & $0.117$ ${[0.055, 0.180]}$ \\
Qwen3-14B & comp., int., seq. & 384 & $4.785$ ${[4.244, 5.354]}$ & 0.167 & 0.169 & $0.003$ ${[-0.039, 0.044]}$ \\
Gemma 4 12B & comp., int., seq. & 384 & $3.949$ ${[3.455, 4.473]}$ & 0.148 & 0.104 & $-0.044$ ${[-0.089, -0.003]}$ \\
Qwen3.5 9B & comp., seq. & 256 & $1.085$ ${[0.975, 1.190]}$ & 0.230 & 0.215 & $-0.016$ ${[-0.078, 0.043]}$ \\
Granite 4.1 8B & int., seq. & 256 & $3.441$ ${[2.977, 3.973]}$ & 0.168 & 0.355 & $0.188$ ${[0.129, 0.246]}$ \\
\bottomrule
\end{tabular}
\caption{Routing results on the confirmation split, using layouts that passed screening on a separate split for each model. $n$ counts paired histories. Seq., int., and comp. denote sequential, interleaved, and competing. S is the superseded assignment and M is the entity mention. Net is the increase in wrong-route frequency with the wrong-route word. Intervals are 97.5\% for the S score shift and 95\% for the paired difference in net switches. For Gemma 4 12B, another interval calculation on the same histories includes zero (see appendix).}
\label{tab:routing}
\end{table*}

Each record contains final routing states for two entities, earlier lines, and updates to a separate archive tag. The policy maps four states through bins to two actions. Table~\ref{tab:example} shows the paired edit. The query, policy, final states, and archive-tag updates are fixed. Within a construction, only the target line's edited word changes. Across constructions, the syntax of both earlier lines changes, and O, M, and U contain no earlier routing-state assignment. Figure~\ref{app:prompt} in the appendix gives the full prompt and both replacement lines for each construction. The wrong-action word always equals the other entity's current routing state, and the correct-action word always equals the word on its earlier line. The entity names differ by one final character. A wrong binary action is therefore compatible with sensitivity to the obsolete word, confusion between entities, or both. The paired contrast measures sensitivity to the edit in the presence of this overlap.

The grid crosses natural and opaque state labels, target action, historical and current order, and entity orientation. Three cells vary record layout: sequential (initial entries before final blocks), interleaved (entries grouped by entity), and competing (interleaved with extra archived discussions quoting state words). State assignments rotate across roles. There are 128 confirmation histories per task and cell, or 384 per full task grid. Each history includes paired items for S, O, M, and U, a current-only item and a pair that edits the live state.

\paragraph{Eligibility and confirmation.}
Development, gate, and confirmation use separate seeds and entity names. Confirmation uses held-out introduction and query wording. The frozen gate requires at least 95\% strict current-state retrieval, at least 90\% current-only action accuracy, and at most 5\% invalid answers in each of direct retrieval and action generation for a cell. A model proceeds if any cell passes. Main cross-model estimates use confirmation histories in the cells selected by that gate. Confirmation competence is assessed separately. Four of the five confirmed models passed the competence checks on different routing layouts in the gate and confirmation splits. Gemma 3 4B and Ministral 3 14B failed every gate cell. OLMo 3 7B stopped at numerical preflight before gate scoring. Five models completed confirmation.

\paragraph{Action scores.}
For prompt $x$, let $L(x)$ be the log probability of the wrong candidate action minus that of the correct candidate. Each event is the action label followed by its model-specific end-of-turn suffix, including a newline where required. For construction $c$, the effect $D_c$ is the history mean of $L(x_{\rm wrong})-L(x_{\rm correct})$, where the subscripts refer to the action implied by the edited word. A positive value means the wrong-action word shifts the relative action score toward the wrong answer, in nats. The original primary outcomes were $D_S$ and $D_S-D_O$. Their 97.5\% intervals adjust for testing both routing and access. Generated switches and mention contrasts are secondary and use 95\% intervals without correction across models. The percentile bootstrap uses 2,000 draws at the history level, keeping the paired members and constructions together. An interval containing zero is inconclusive.

\paragraph{Generated decisions.}
Greedy generation uses a 16-token budget with thinking disabled in the main runs. After whitespace stripping, outputs must equal an allowed label. Define net switch rate as wrong-action frequency with the wrong-action word minus wrong-action frequency with the correct-action word. Equivalently, count transitions into the wrong-action category minus transitions out, then divide by histories. Invalid answers are outside the wrong-action category in this calculation and count as failures for strict accuracy and prediction. Transitions involving invalid answers can therefore contribute to the metric. It measures a net change and can conceal switches in both directions.

\paragraph{Separate retrieval questions.}
We ask directly for the current state of each entity and separately for the initial state in the superseded prompts. For a post hoc tracking check, require correct target-current answers for both members of each included construction, correct initial answers for both superseded members, and a correct current-only action. Within each model, compare constructions on the intersection of histories passing these checks under all four constructions.

\section{Results}
Editing a superseded assignment shifted action scores and increased the frequency of wrong decisions in all five models on the gate-selected routing histories (Table~\ref{tab:routing}). Editing an incidental mention also increased wrong decisions in all five, with 95\% intervals above zero. For example, in Qwen3-8B the increase was about 20 percentage points for the superseded assignment and 31 for the mention, on the same 256 histories. Mention edits also shifted scores toward the wrong action in each model on the full routing grid.

\paragraph{The difference between mention and assignment varies.}
On the gate-selected histories, Qwen3-8B and Granite 4.1 8B had larger net effects for mentions, and the paired differences were inconclusive for Qwen3-14B and Qwen3.5 9B. Gemma 4 12B had a smaller estimate for mentions, with the interval's exclusion of zero sensitive to the analysis (see appendix).

\begin{table*}[t]
\centering
\begin{tabular}{@{}lrrrrrrrr@{}}
\toprule
Model & S net & O net & M net & U net & \shortstack{S wrong\\word} & \shortstack{S correct\\word} & \shortstack{M wrong\\word} & \shortstack{Current\\only} \\
\midrule
Qwen3-8B & 0.185 & 0.065 & 0.294 & 0.039 & 74.0 & 92.4 & 64.3 & 95.6 \\
Qwen3-14B & 0.167 & 0.146 & 0.169 & 0.029 & 72.1 & 88.8 & 73.4 & 98.4 \\
Gemma 4 12B & 0.148 & 0.010$^{*}$ & 0.104 & 0.016$^{*}$ & 75.5 & 90.4 & 83.3 & 90.9 \\
Qwen3.5 9B & 0.229 & 0.109 & 0.193 & 0.026$^{*}$ & 64.1 & 87.0 & 71.4 & 93.5 \\
Granite 4.1 8B & 0.164 & 0.138 & 0.370 & 0.182 & 69.3 & 85.7 & 52.6 & 91.1 \\
\bottomrule
\end{tabular}
\caption{Full routing grid, 384 confirmation histories per model. Net switches are fractions. The last four columns give strict action accuracy in percent. An asterisk marks a net rate whose 95\% paired history-bootstrap interval includes zero. All use the same histories, including layouts outside the frozen gate selection.}
\label{tab:full_grid}
\end{table*}

\paragraph{Full-grid accuracy and controls.}
Table~\ref{tab:full_grid} includes every routing layout. On this grid, 75 of 384 Qwen3-8B superseded pairs switched toward the wrong action and four switched away, a net rate of 0.185 [0.143, 0.227]. For Granite 4.1 8B, even the unattached edit has a net switch rate of 0.182, compared with 0.164 for S. These are descriptive marginal estimates. For all five models, strict accuracy is lower with the wrong-action word than with the correct-action word in S. The current-only item benchmarks the action on a resolved record. It removes both earlier lines and, in competing layouts, the archived discussions. A matched neutral-word member would be needed to distinguish harm from the wrong-action word from help from the correct-action word.

\paragraph{Correct recall can coexist with action changes.}
On each model's post hoc common tracked set, the superseded edit still shows positive net switching. For example, Qwen3-8B retains 335 of 384 histories, with net rates of 0.161 for superseded assignments and 0.266 for incidental mentions (paired difference 0.104 [0.057, 0.155]). The appendix gives the other retained counts and comparisons. The model could give both values when asked directly. Those answers tell us little about which value it used to choose the action.

\paragraph{A never-updated attribute also interferes.}
The primary score contrast favors S over O in all five routing models. Both constructions include an update to the attribute that held the edited word, and the update to O could itself reduce the word's influence. A follow-up on the full 384 routing histories per Qwen3 model leaves the other attribute unchanged (N), with the direction of $D_O-D_N$ recorded before scoring. Updating the other attribute reduced its score effect by 1.204 nats [1.080, 1.326] for Qwen3-8B and 1.942 [1.706, 2.191] for Qwen3-14B (95\% intervals). The paired $D_S-D_N$ was inconclusive for Qwen3-8B (0.020 [$-0.201$, 0.254]) and favored N for Qwen3-14B ($-1.076$ [$-1.526$, $-0.661$]). The S--O contrast alone is therefore insufficient to establish specificity to the queried attribute.

\paragraph{Thinking-enabled generation.}
On the original 384 routing histories per model, exploratory thinking-enabled runs observed zero switches in either direction for S in both Qwen3 models. M still had seven wrong-directed switches in Qwen3-8B and one in Qwen3-14B, with zero reverse switches. The generation runs returned two invalid answers out of 3,456 items in Qwen3-8B and zero in Qwen3-14B. We used a longer generation budget and recorded answers only. Whether thinking reduces switches on other records remains open.

\section{Implications and limitations}
A regression test for a system that reads versioned records can start with a case whose correct action is known. Keep the final record and policy fixed. Change an old value, then try the same edit in a note that mentions it. Count how often decisions get worse or better and record invalid answers. Ask for the current state separately. Use the same record format and decoding settings the deployed system will see.

The results apply to short synthetic records and the configurations tested here. Each construction uses a fixed sentence frame, and models can pass the gate on different layouts. In Qwen3-8B, the mention score effect exceeded the superseded score effect on confirmation, while the gate comparison was inconclusive. Both the wording and the histories differed between splits. A separate follow-up found fewer superseded switches on new histories where both entities had the same current state. These comparisons leave open how much the results depend on wording, entity overlap, and task difficulty.

\bibliography{references}
\appendix
\section{Additional conditions and competence}
Table~\ref{tab:gate} gives the frozen gate checks. Table~\ref{tab:confirmation} applies the competence criteria to the held-out confirmation split. A change in competence between splits leaves the gate-selected population fixed. The original design used separate development, gate, and confirmation splits with 192, 384, and 768 histories, respectively, across both tasks.

\paragraph{Access task.}
Table~\ref{tab:access} gives full-grid estimates for the five confirmed models. In the access task the final clearance maps directly to APPROVE or DENY. Only Gemma 4 12B met confirmation competence in access, passing all three cells. Its gate-selected competing cell (128 histories) gave $D_S=4.384$ [3.480, 5.518] nats and net switching of 0.062 [0.023, 0.109]. Qwen3-8B passed the frozen gate in the interleaved cell (128 histories) but failed confirmation competence there. Its corresponding estimates were 4.221 [3.470, 5.063] nats and 0.273 [0.195, 0.352]. The score intervals are 97.5\%, and switch intervals are 95\%. All other confirmed models failed the access gate in every cell. The confirmation action outputs included one invalid item in the Granite 4.1 8B competing access cell and zero invalid items in routing for all five models.

\paragraph{Balanced histories.}
Table~\ref{tab:conditions} reports all four constructions. An exploratory amendment introduced 192 new routing histories for Qwen3-8B, with identical current states for the target and other entity. The S score shift was 3.895 [3.333, 4.443] nats (95\% interval), and six pairs switched into the wrong-action category and zero out, for a net of 0.031 [0.010, 0.057]. M produced 34 wrong-directed switches and zero reverse switches, for a net of 0.177 [0.125, 0.234]. This was a descriptive follow-up with no prespecified test.

\paragraph{Prediction analysis.}
The secondary H3 analysis asks whether a retrieval-score sensitivity term improves prediction of strict downstream failure. For each entity, measure the paired change in the log-mass margin between the two edited words on the direct current-state query. Subtract the other entity's sensitivity from the target's. Retrieval scoring sums bounded answer-prefix masses across surface forms. The baseline uses pre-decision retrieval scores, word position, prompt length, and design indicators. Both models use fixed L2-regularized logistic regression and five history-grouped folds, with the paired difference in held-out log loss bootstrapped by history. Table~\ref{tab:h3} reports each task separately. Routing gains are inconclusive under this procedure.

\section{Reproducibility}
The confirmation protocol was frozen internally before confirmation. Its revision history includes numerical and answer-event formatting corrections before the final gate. Scoring used unquantized float32 with SDPA on one 80 GB H100, with cached action scores checked against full forwards. Generation-only thinking runs used up to 4,096 new tokens and parsed the answer after the last closing thinking marker. Unclosed, truncated, and unparseable outputs counted as invalid. The balanced follow-up used the confirmation runtime and accepted one small causality deviation under its runner's tolerance. OLMo 3 7B was excluded under a preflight that required exactly zero deviation.

The original analysis used saved score files, manifests, and completion records with hash checks. The tables here use saved derived summaries and their paired intervals. Regenerating the tables from those summaries checks the reporting layer. The sealed input files were unavailable for this build, so raw responses and score provenance could not be independently verified. An anonymous artifact link is omitted.

\begin{table*}[t]
\centering
\begin{tabular}{@{}lcccccc@{}}
\toprule
Role & Main score & Main net & Bal. score & Bal. net & Think 8B & Think 14B \\
\midrule
S & \shortstack{$3.133$\\$[2.794, 3.478]$} & \shortstack{$0.185$\\$[0.143, 0.227]$} & \shortstack{$3.895$\\$[3.333, 4.443]$} & \shortstack{$0.031$\\$[0.010, 0.057]$} & 0/0 & 0/0 \\
O & \shortstack{$1.909$\\$[1.661, 2.161]$} & \shortstack{$0.065$\\$[0.029, 0.099]$} & \shortstack{$2.962$\\$[2.541, 3.365]$} & \shortstack{$0.026$\\$[0.005, 0.052]$} & 0/0 & 2/0 \\
M & \shortstack{$4.807$\\$[4.455, 5.186]$} & \shortstack{$0.294$\\$[0.250, 0.341]$} & \shortstack{$6.621$\\$[5.917, 7.323]$} & \shortstack{$0.177$\\$[0.125, 0.234]$} & 7/0 & 1/0 \\
U & \shortstack{$2.496$\\$[2.201, 2.805]$} & \shortstack{$0.039$\\$[0.013, 0.068]$} & \shortstack{$3.326$\\$[2.889, 3.765]$} & \shortstack{$-0.005$\\$[-0.026, 0.010]$} & 0/0 & 0/0 \\
\bottomrule
\end{tabular}
\caption{Full-grid routing conditions. Main (384 histories) and balanced (192) use Qwen3-8B without thinking. Scores are in nats. The displayed intervals are secondary 95\% intervals. Thinking shows counts into/out of the wrong-action category on 384 main histories per model, with no action scores.}
\label{tab:conditions}
\end{table*}
\begin{table*}[t]
\centering
\begin{tabular}{@{}llcc@{}}
\toprule
Model & Competent cells & $D_S$ (nats) & S net \\
\midrule
Qwen3-8B & none & $3.904$ ${[3.519, 4.306]}$ & $0.167$ ${[0.130, 0.206]}$ \\
Qwen3-14B & none & $4.377$ ${[3.806, 4.960]}$ & $0.161$ ${[0.125, 0.201]}$ \\
Gemma 4 12B & all three & $4.285$ ${[3.699, 4.896]}$ & $0.083$ ${[0.057, 0.112]}$ \\
Qwen3.5 9B & none & $1.444$ ${[1.323, 1.573]}$ & $0.227$ ${[0.185, 0.271]}$ \\
Granite 4.1 8B & none & $3.294$ ${[2.772, 3.840]}$ & $0.068$ ${[0.042, 0.096]}$ \\
\bottomrule
\end{tabular}
\caption{Full-grid access confirmation (384 histories per model, including cells failing competence). Both outcomes use 95\% intervals for these descriptive construction summaries. The gate-selected primary access intervals are reported in the text.}
\label{tab:access}
\end{table*}
\begin{table*}[t]
\centering
\begin{tabular}{@{}lrccccc@{}}
\toprule
Model & Retained & S & O & M & U & M $-$ S \\
\midrule
Qwen3-8B & 335/384 & 0.161 & 0.045 & 0.266 & 0.033 & $0.104$ ${[0.057, 0.155]}$ \\
Qwen3-14B & 359/384 & 0.170 & 0.139 & 0.167 & 0.025 & $-0.003$ ${[-0.050, 0.042]}$ \\
Gemma 4 12B & 339/384 & 0.147 & 0.006 & 0.068 & 0.009 & $-0.080$ ${[-0.124, -0.038]}$ \\
Qwen3.5 9B & 294/384 & 0.228 & 0.085 & 0.177 & 0.044 & $-0.051$ ${[-0.102, 0.000]}$ \\
Granite 4.1 8B & 299/384 & 0.157 & 0.127 & 0.321 & 0.184 & $0.164$ ${[0.117, 0.217]}$ \\
\bottomrule
\end{tabular}
\caption{Net switch rates on the post hoc common tracked routing histories within each model. Paired differences have 95\% intervals.}
\label{tab:tracked}
\end{table*}
\begin{table*}[t]
\centering
\begin{tabular}{@{}lllrrrrrc@{}}
\toprule
Model & Task & Cell & Direct & Current-only & History & Invalid & Easy invalid & Pass \\
\midrule
Qwen3-8B & routing & seq. & 0.977 & 0.938 & 0.779 & 0.000 & 0.000 & yes \\
Qwen3-8B & routing & int. & 0.986 & 0.938 & 0.779 & 0.000 & 0.001 & yes \\
Qwen3-8B & routing & comp. & 0.945 & 0.938 & 0.752 & 0.000 & 0.000 & no \\
Qwen3-8B & access & seq. & 0.889 & 0.906 & 0.822 & 0.000 & 0.055 & no \\
Qwen3-8B & access & int. & 0.953 & 0.938 & 0.877 & 0.000 & 0.013 & yes \\
Qwen3-8B & access & comp. & 0.953 & 0.859 & 0.883 & 0.000 & 0.009 & no \\
Qwen3-14B & routing & seq. & 0.996 & 0.922 & 0.854 & 0.000 & 0.000 & yes \\
Qwen3-14B & routing & int. & 1.000 & 1.000 & 0.846 & 0.000 & 0.000 & yes \\
Qwen3-14B & routing & comp. & 0.997 & 1.000 & 0.814 & 0.000 & 0.000 & yes \\
Qwen3-14B & access & seq. & 0.991 & 0.781 & 0.729 & 0.000 & 0.001 & no \\
Qwen3-14B & access & int. & 0.977 & 0.828 & 0.748 & 0.000 & 0.001 & no \\
Qwen3-14B & access & comp. & 0.946 & 0.828 & 0.697 & 0.000 & 0.000 & no \\
Gemma 3 4B & routing & seq. & 0.939 & 0.688 & 0.645 & 0.000 & 0.023 & no \\
Gemma 3 4B & routing & int. & 0.964 & 0.609 & 0.488 & 0.000 & 0.003 & no \\
Gemma 3 4B & routing & comp. & 0.916 & 0.625 & 0.533 & 0.000 & 0.001 & no \\
Gemma 3 4B & access & seq. & 0.945 & 0.719 & 0.678 & 0.000 & 0.027 & no \\
Gemma 3 4B & access & int. & 0.936 & 0.609 & 0.676 & 0.000 & 0.021 & no \\
Gemma 3 4B & access & comp. & 0.946 & 0.672 & 0.641 & 0.000 & 0.021 & no \\
Ministral 3 14B & routing & seq. & 0.997 & 0.766 & 0.766 & 0.000 & 0.000 & no \\
Ministral 3 14B & routing & int. & 0.986 & 0.766 & 0.756 & 0.000 & 0.000 & no \\
Ministral 3 14B & routing & comp. & 0.999 & 0.812 & 0.756 & 0.000 & 0.000 & no \\
Ministral 3 14B & access & seq. & 0.970 & 0.750 & 0.670 & 0.000 & 0.000 & no \\
Ministral 3 14B & access & int. & 0.926 & 0.812 & 0.711 & 0.000 & 0.003 & no \\
Ministral 3 14B & access & comp. & 0.952 & 0.797 & 0.732 & 0.000 & 0.006 & no \\
Gemma 4 12B & routing & seq. & 1.000 & 0.906 & 0.936 & 0.000 & 0.000 & yes \\
Gemma 4 12B & routing & int. & 0.999 & 0.984 & 0.887 & 0.000 & 0.000 & yes \\
Gemma 4 12B & routing & comp. & 1.000 & 0.969 & 0.930 & 0.000 & 0.000 & yes \\
Gemma 4 12B & access & seq. & 0.912 & 0.969 & 0.967 & 0.000 & 0.088 & no \\
Gemma 4 12B & access & int. & 0.987 & 0.875 & 0.850 & 0.000 & 0.011 & no \\
Gemma 4 12B & access & comp. & 0.993 & 0.906 & 0.932 & 0.000 & 0.003 & yes \\
Qwen3.5 9B & routing & seq. & 1.000 & 0.906 & 0.795 & 0.000 & 0.000 & yes \\
Qwen3.5 9B & routing & int. & 0.922 & 0.922 & 0.746 & 0.000 & 0.000 & no \\
Qwen3.5 9B & routing & comp. & 0.950 & 0.969 & 0.867 & 0.000 & 0.000 & yes \\
Qwen3.5 9B & access & seq. & 0.994 & 0.688 & 0.656 & 0.000 & 0.000 & no \\
Qwen3.5 9B & access & int. & 0.908 & 0.703 & 0.666 & 0.000 & 0.000 & no \\
Qwen3.5 9B & access & comp. & 0.891 & 0.688 & 0.789 & 0.000 & 0.000 & no \\
Granite 4.1 8B & routing & seq. & 0.967 & 0.984 & 0.799 & 0.000 & 0.000 & yes \\
Granite 4.1 8B & routing & int. & 0.955 & 0.953 & 0.738 & 0.000 & 0.000 & yes \\
Granite 4.1 8B & routing & comp. & 0.932 & 0.953 & 0.816 & 0.000 & 0.000 & no \\
Granite 4.1 8B & access & seq. & 0.935 & 0.797 & 0.828 & 0.000 & 0.000 & no \\
Granite 4.1 8B & access & int. & 0.960 & 0.812 & 0.793 & 0.000 & 0.000 & no \\
Granite 4.1 8B & access & comp. & 0.962 & 0.828 & 0.865 & 0.000 & 0.000 & no \\
\bottomrule
\end{tabular}
\caption{Gate competence. Each cell contains 64 histories. Direct is target-current retrieval accuracy. Easy invalid is the fraction of invalid retrieval answers. Both use 11 items per history, as does the action invalid rate. Current-only accuracy uses one action item and historical accuracy uses eight.}
\label{tab:gate}
\end{table*}
\begin{table*}[t]
\centering
\begin{tabular}{@{}lllrrrrrc@{}}
\toprule
Model & Task & Cell & Direct & Current-only & History & Invalid & Easy invalid & Pass \\
\midrule
Qwen3-8B & routing & seq. & 0.988 & 0.953 & 0.858 & 0.000 & 0.000 & yes \\
Qwen3-8B & routing & int. & 0.986 & 0.961 & 0.868 & 0.000 & 0.000 & yes \\
Qwen3-8B & routing & comp. & 0.982 & 0.953 & 0.846 & 0.000 & 0.000 & yes \\
Qwen3-8B & access & seq. & 0.916 & 0.867 & 0.831 & 0.000 & 0.049 & no \\
Qwen3-8B & access & int. & 0.945 & 0.875 & 0.858 & 0.000 & 0.026 & no \\
Qwen3-8B & access & comp. & 0.908 & 0.859 & 0.867 & 0.000 & 0.049 & no \\
Qwen3-14B & routing & seq. & 0.989 & 0.984 & 0.871 & 0.000 & 0.000 & yes \\
Qwen3-14B & routing & int. & 0.996 & 0.977 & 0.851 & 0.000 & 0.000 & yes \\
Qwen3-14B & routing & comp. & 0.997 & 0.992 & 0.843 & 0.000 & 0.000 & yes \\
Qwen3-14B & access & seq. & 0.919 & 0.836 & 0.744 & 0.000 & 0.055 & no \\
Qwen3-14B & access & int. & 0.920 & 0.859 & 0.776 & 0.000 & 0.049 & no \\
Qwen3-14B & access & comp. & 0.962 & 0.820 & 0.747 & 0.000 & 0.024 & no \\
Gemma 4 12B & routing & seq. & 0.999 & 0.914 & 0.898 & 0.000 & 0.001 & yes \\
Gemma 4 12B & routing & int. & 0.997 & 0.898 & 0.840 & 0.000 & 0.001 & no \\
Gemma 4 12B & routing & comp. & 0.977 & 0.914 & 0.905 & 0.000 & 0.022 & yes \\
Gemma 4 12B & access & seq. & 0.969 & 0.969 & 0.937 & 0.000 & 0.030 & yes \\
Gemma 4 12B & access & int. & 0.988 & 0.992 & 0.929 & 0.000 & 0.010 & yes \\
Gemma 4 12B & access & comp. & 0.978 & 0.969 & 0.934 & 0.000 & 0.007 & yes \\
Qwen3.5 9B & routing & seq. & 0.986 & 0.930 & 0.770 & 0.000 & 0.000 & yes \\
Qwen3.5 9B & routing & int. & 0.954 & 0.906 & 0.836 & 0.000 & 0.000 & yes \\
Qwen3.5 9B & routing & comp. & 0.923 & 0.969 & 0.877 & 0.000 & 0.003 & no \\
Qwen3.5 9B & access & seq. & 0.998 & 0.805 & 0.748 & 0.000 & 0.002 & no \\
Qwen3.5 9B & access & int. & 0.997 & 0.789 & 0.759 & 0.000 & 0.000 & no \\
Qwen3.5 9B & access & comp. & 0.965 & 0.805 & 0.846 & 0.000 & 0.009 & no \\
Granite 4.1 8B & routing & seq. & 0.955 & 0.898 & 0.796 & 0.000 & 0.000 & no \\
Granite 4.1 8B & routing & int. & 0.957 & 0.906 & 0.744 & 0.000 & 0.000 & yes \\
Granite 4.1 8B & routing & comp. & 0.963 & 0.930 & 0.790 & 0.000 & 0.000 & yes \\
Granite 4.1 8B & access & seq. & 0.932 & 0.883 & 0.866 & 0.000 & 0.000 & no \\
Granite 4.1 8B & access & int. & 0.920 & 0.828 & 0.824 & 0.000 & 0.000 & no \\
Granite 4.1 8B & access & comp. & 0.949 & 0.867 & 0.893 & 0.001 & 0.000 & no \\
\bottomrule
\end{tabular}
\caption{Confirmation competence. Each cell contains 128 histories. Direct is target-current retrieval accuracy. Easy invalid is the fraction of invalid retrieval answers. Both use 11 items per history, as does the action invalid rate. Current-only accuracy uses one action item and historical accuracy uses eight.}
\label{tab:confirmation}
\end{table*}
\begin{table*}[t]
\centering
\begin{tabular}{@{}llrc@{}}
\toprule
Model & Task & Items & Log-loss gain \\
\midrule
Qwen3-8B & routing & 3072 & $0.0013$ ${[-0.0022, 0.0047]}$ \\
Qwen3-8B & access & 3072 & $-0.0006$ ${[-0.0020, 0.0008]}$ \\
Qwen3-14B & routing & 3072 & $0.0105$ ${[0.0022, 0.0184]}$ \\
Qwen3-14B & access & 3072 & $-0.0001$ ${[-0.0023, 0.0020]}$ \\
Gemma 4 12B & routing & 3072 & $0.0008$ ${[-0.0028, 0.0041]}$ \\
Gemma 4 12B & access & 3072 & $-0.0005$ ${[-0.0024, 0.0012]}$ \\
Qwen3.5 9B & routing & 3072 & $-0.0003$ ${[-0.0025, 0.0019]}$ \\
Qwen3.5 9B & access & 3072 & $0.0038$ ${[-0.0011, 0.0086]}$ \\
Granite 4.1 8B & routing & 3072 & $-0.0000$ ${[-0.0013, 0.0013]}$ \\
Granite 4.1 8B & access & 3072 & $-0.0025$ ${[-0.0050, -0.0004]}$ \\
\bottomrule
\end{tabular}
\caption{H3 on the full confirmation grid (384 histories per model and task). Positive gain means lower out-of-fold log loss with retrieval sensitivity. Paired 95\% history-bootstrap intervals.}
\label{tab:h3}
\end{table*}


\end{document}
